\documentclass{article}
\usepackage{blindtext}
\usepackage[a4paper, total={6in, 8in}]{geometry}
\usepackage{graphicx}
\usepackage{colortbl, color, xcolor}
\usepackage{hyperref}
\usepackage{cleveref}
\usepackage{booktabs}
\usepackage{enumitem}

\title{NTIRE2026 -- Robust Deepfake Detection Challenge}

\author{[YOUR TEAM NAME HERE]\\
[YOUR NAME]\thanks{\textcolor{blue}{[YOUR AFFILIATION]}}\\
}

\begin{document}
\maketitle

\section{Team details}

\subsection*{[YOUR TEAM NAME]}
\noindent\textit{\textbf{Title: }} Frequency-Aware Multi-Branch Deepfake Detector with Patch-Based Training\\
\noindent\textit{\textbf{Members:}}\\ 
\textit{[TEAM LEADER]$^1$(\href{mailto:[EMAIL]}{[EMAIL]})}, [OTHER MEMBERS]$^{1}$\\
\noindent\textit{\textbf{Affiliations: }}\\
$^1$ [YOUR INSTITUTION/AFFILIATION]\\

\noindent\textit{\textbf{CodaBench Username:}} [YOUR USERNAME]\\
\noindent\textit{\textbf{Code Repository:}} [GITHUB/DRIVE LINK]

\section{Contribution details}

\subsection{Overall Approach}
Our solution combines three complementary detection strategies to achieve robustness against image degradations:

\textbf{1. Spatial Feature Extraction:} ConvNeXt-Tiny backbone with attention-based pooling extracts texture and artifact patterns from the spatial domain.

\textbf{2. Frequency Analysis:} Fast Fourier Transform (FFT) branch analyzes magnitude spectra to detect GAN upsampling artifacts that persist in frequency domain even under heavy spatial degradation (noise, blur, compression).

\textbf{3. Patch-Based Training:} During training, we randomly extract 4 patches (128×128) per image and average predictions. This prevents the model from memorizing global face identity and forces it to learn local manipulation artifacts.

\textbf{Key Innovation:} The combination of FFT frequency analysis with patch-based training addresses the core challenge — maintaining detection performance when images are degraded by noise, blur, low resolution, and JPEG compression.

\subsection{Architecture Details}
\begin{itemize}
    \item \textbf{Backbone:} ConvNeXt-Tiny (28.8M parameters) with ImageNet-1K pre-training
    \item \textbf{FFT Branch:} 2D Fourier transform → magnitude spectrum → 3-layer CNN encoder (256-dim output)
    \item \textbf{Fusion:} Concatenate spatial + frequency features → 3-layer MLP with LayerNorm and dropout
    \item \textbf{Training Strategy:} 
    \begin{itemize}
        \item Random patch extraction (4 patches/image, 128×128)
        \item Label smoothing (0.1) to prevent overconfidence
        \item Strong augmentation: Gaussian noise, blur, brightness/contrast, horizontal flip
        \item 85/15 train/val split with cosine LR schedule
    \end{itemize}
\end{itemize}

\subsection{References}
\begin{itemize}
    \item ConvNeXt: Liu et al., "A ConvNet for the 2020s", CVPR 2022
    \item Frequency analysis for deepfakes: Frank et al., "Leveraging Frequency Analysis for Deep Fake Image Recognition", ICML 2020
    \item Patch-based robustness: Inspired by masked image modeling approaches
\end{itemize}

\section{Global Method Facts}

\begin{enumerate}[label=\alph*)]
\item \textbf{Total complexity:} 28.8M parameters (ConvNeXt) + 0.5M (FFT branch + fusion head) = 29.3M total
\item \textbf{Pre-trained models:} ConvNeXt-Tiny pre-trained on ImageNet-1K from timm library
\item \textbf{Extra data:} No external datasets used beyond the provided 1000 training images
\item \textbf{Training:} 
\begin{itemize}
    \item 30 epochs with early stopping (patience=10)
    \item Batch size 32, effective batch 32 (no gradient accumulation)
    \item AdamW optimizer: LR=5e-5, weight decay=1e-4
    \item Cosine annealing schedule with warmup
    \item Training time: ~30 minutes on single GPU
\end{itemize}
\item \textbf{Testing:} 
\begin{itemize}
    \item 3-view TTA: original + horizontal flip + center crop from 1.1× scale
    \item Inference time: ~2 minutes for 100 images
    \item No patch extraction at test time (full image used)
\end{itemize}
\item \textbf{Advantages:}
\begin{itemize}
    \item Frequency branch maintains detection under heavy spatial degradation
    \item Patch-based training prevents identity memorization and improves generalization
    \item Lightweight (29M params) enables fast training and inference
    \item No dependence on external datasets
\end{itemize}
\item \textbf{Novelty:} Combination of FFT frequency analysis with patch-based training for robustness. While each component exists in literature, their integration for degradation-robust deepfake detection is novel.
\end{enumerate}

\begin{table}[htb]
    \centering
    \resizebox{\textwidth}{!}{
    \begin{tabular}{c|c|c|c|c|c|c|c|c|c}
        Input & Training Time & Epochs & Extra data & Attention & Quantization & Base Model & \# Params. (M) & Runtime & GPU  \\
        \hline
        (224, 224, 3) & 30 min & 30 & No & No & No & ConvNeXt-Tiny & 29.3M & 1.2s (100 imgs) & RTX 3060/4060 \\
    \end{tabular}
    }
    \caption{Method specification}
    \label{tab:methodinfos}
\end{table}

\section{Technical details}

\textbf{Implementation:}
\begin{itemize}
    \item Framework: PyTorch 2.0+
    \item Libraries: timm (models), albumentations (augmentation), scikit-learn (metrics)
    \item Optimizer: AdamW with cosine annealing LR schedule
    \item Loss: Binary cross-entropy with label smoothing (0.1)
    \item GPU: Trained on NVIDIA RTX 3060/4060 (12GB VRAM sufficient)
    \item Dataset: 1000 images (500 real, 500 fake) split 85/15 train/val
\end{itemize}

\textbf{Training strategy:}
\begin{itemize}
    \item Patch extraction during training only (4 random 128×128 patches per image)
    \item Strong augmentation to simulate test-time degradations
    \item Early stopping based on validation AUC (patience=10 epochs)
    \item No data augmentation mixup or cutmix (patch extraction provides sufficient regularization)
\end{itemize}

\section{Other details}

\begin{enumerate}[label=\alph*)]
\item \textbf{Workshop paper submission:} [YES/NO] - Please update

\item \textbf{Challenge feedback:} The challenge effectively tests robustness to real-world degradations. The validation/test split with different degradation levels appropriately assesses generalization. Suggestion: provide degradation type labels on test set for post-hoc analysis.

\item \textbf{Future challenge suggestions:} 
\begin{itemize}
    \item Video-based deepfake detection with temporal artifacts
    \item Multi-modal detection (image + audio)
    \item Adversarial robustness track (detection under adversarial attacks)
\end{itemize}
\end{enumerate}

\end{document}
